% Chapter 11 -- Alerts and Dialog Boxes (Part II, Reference; the FLAIR
% counterpart of a messages chapter).  One \AlertEntry per alert or dialog
% box the desktop and its applications show, keyed by its first line as the
% source draws it; the keys are compared with harness/pubs/flair_alerts.sh
% (the source's own literals) by test-pubs-flair in both directions.
% Sources (Law 1): os/milton/kmain.c (finder_trash_confirm: the Locked
% notice and the Empty Trash buttons; finder_shutdown_screen; the About
% alert flair_live_about with BANNER_LINE1; finder_exit_dos with Exit /
% Cancel; sf_run's Replace question with Replace / Cancel);
% os/flair/finder_windows.c (finder_trash_alert_text, the singular and
% plural wording, K rounded up; finder_exit_dos_alert_text);
% os/flair/dialog.h (FLAIR_CANON_FILECOPY_TITLE / _MSG);
% os/apps/itext/itext.c (say(): the three notices; guarded(): the Save
% changes question, "before closing?" / "before quitting?"; the Save
% prompt).  Every text quoted character for character (NT-16).
% GATES: test-flair-empty-trash test-flair-locked-trash test-flair-shutdown test-flair-apple-menu test-flair-exit-to-dos test-stdfile test-flair-itext test-itext-route test-flair-file-copy test-pubs-flair test-pubs-flair-screens
\Gates{test-flair-empty-trash, test-flair-locked-trash, test-flair-shutdown,
  test-flair-apple-menu, test-flair-exit-to-dos, test-stdfile,
  test-flair-itext, test-itext-route, test-flair-file-copy, test-pubs-flair,
  test-pubs-flair-screens}
\chapter{Alerts and Dialog Boxes}

This chapter lists, in alphabetical order, every alert, dialog box and
progress box that FLAIR and the applications described in this book and the
\doctitle{FLAIR Accessories Guide} display.  Each is shown by its text,
exactly as it appears on the screen; words in italic stand for what FLAIR
fills in.  After each comes what it means and what to do.

In every alert, the button with the heavy border is chosen by \key{Enter},
and \texttt{Cancel}, where there is one, by \key{Esc}.

\AlertEntry{Are you sure you want to exit to InitechDOS?}{}{Shown when you
choose \menuitem{Exit to InitechDOS} from the Finder's \menuname{Special}
menu.  Buttons \texttt{Exit} and \texttt{Cancel}.  \texttt{Exit} asks an
open application to quit, saves the arrangement of the desktop and leaves
FLAIR for the InitechDOS command prompt; \texttt{Cancel} leaves everything
as it was.  See Chapter~1.}

\AlertEntry{FILE COPY}{FILE COPY}{The title of the progress box InitechCalc
shows while it saves its tables (see \texttt{Saving tables to disk...}).}

\AlertEntry{InitechDOS Version 3.30}{InitechDOS\ \ Version 3.30\\
Built-in Memory: \textit{n}K\\
Largest Unused Block: \textit{n}K}{The About This Computer alert, shown when
you choose \menuitem{About This Computer...} from the Apple menu.  The first
line is the version of InitechDOS on the startup disk.  \textit{Built-in
Memory} is all the memory of the computer, in kilobytes; \textit{Largest
Unused Block} is the largest part of the memory InitechDOS gives to
programs that is free at the moment (\texttt{none} if InitechDOS has not
yet given out its memory).  Click \texttt{OK}.}

\AlertEntry{It is now safe to switch off your computer.}{It is now safe
to switch off your computer.\\ If you wish to continue using your
computer,\\ press Restart.}{Shown on a black screen after \menuitem{Shut
Down}.  Turn off the computer, or click \texttt{Restart} (or press
\key{Enter}) to start it again.}

\AlertEntry{Replace existing}{Replace existing "\textit{NAME}"?}{Shown by
the Save dialog box when the name you gave is already used by a file in
the folder.  \texttt{Replace} writes the document over that file, which is
lost; \texttt{Cancel} lets you type another name.}

\AlertEntry{Save changes to the document}{Save changes to the document\\
"\textit{name}" before closing?}{Shown by InitechText when you click the
close box of its window or choose \menuitem{New} or \menuitem{Open...},
and --- with the second line ending \texttt{before quitting?} --- when you
choose \menuitem{Close} or \menuitem{Quit}, while the document has changes
you have not saved.  \texttt{Save} saves
them first, \texttt{Don't Save} goes on without saving, and \texttt{Cancel}
returns to the document.}

\AlertEntry{Save this document as:}{}{The prompt above the name box in the
Save dialog box (Chapter~5).}

\AlertEntry{Saving tables to disk...}{}{The message in InitechCalc's
\texttt{FILE COPY} progress box, shown while it copies its ledger to the
disk.  Wait: the box goes away by itself when the copy is finished, and
clicks and keys are ignored until then.  See the \doctitle{FLAIR Accessories
Guide}.}

\AlertEntry{The document could not be opened.}{}{InitechText could not read
the file you chose.  Click \texttt{OK}.  Check that the file is a text file
and that the disk is in place.}

\AlertEntry{The document could not be saved.}{}{InitechText could not write
the file, usually because the disk is full or the file is locked.  Click
\texttt{OK}, then save the document under another name, or make room on the
disk.}

\AlertEntry{The item is locked.}{The item is locked.\\ Unlock it before\\
discarding it.}{You dragged a locked file to the Trash, or chose
\menuitem{Empty Trash} while the Trash held a locked item.  The item is
not thrown away.  Click \texttt{OK}.  To throw the item away, unlock it
first at the InitechDOS prompt with \file{ATTRIB} (Chapter~6).}

\AlertEntry{The Trash contains}{The Trash contains \textit{n} items, which
use \textit{k}K of disk space. Are you sure you want to remove these items
permanently?}{Shown when you choose \menuitem{Empty Trash}.  When the Trash
holds one item, the alert reads \texttt{The Trash contains 1 item, which
uses \textit{k}K of disk space. Are you sure you want to remove this item
permanently?}  \textit{k} is the space the items use, in kilobytes,
rounded up.  \texttt{OK} deletes everything in the Trash; \texttt{Cancel}
keeps it.}

\AlertEntry{This document is too large to be}{This document is too large
to be\\ opened by InitechText.}{The file you chose is longer than
InitechText can hold.  Click \texttt{OK}.  Open the file with an InitechDOS
program instead, or divide it into smaller files.}
